% ---
% id: receiver/component-catalogue/thm:catalogue-admissibility
% kind: thm
% title: "Admissibility of the catalogue entries"
% summary: >-
%   Every catalogue entry maps its whole declared parameter domain to
%   contractions on its modes, lossless exactly on the subfamily the
%   catalogue's table lists for it. Only the generic block, an arbitrary matrix
%   on declared modes, needs its contractivity checked at admission.
% prominence: main
% status: implemented
% uses: [receiver/thm:compilation-well-defined, receiver/def:compile, receiver/def:embed, receiver/def:graph, lem:attenuated, def:catalogue-entries]
% mentions: [rem:bsconv]
% assumes: []
% refs: []
% provenance: newly-derived
% implements: certiqs_key.receiver.catalogue:Component
% lean:
% ---

A \tref{declared element graph} is \tref{well formed} only if every element
is a contraction on its support. For entries C1--C9 this holds by theorem on
the whole declared domain, so for each of them condition (E1) of
Definition~\ref{def:graph} reduces to checking that its parameters lie in
that domain.

\begin{theorem}\label{thm:catalogue-admissibility}
Let $\Gmat_\Sigma$ be the matrix of an entry type C1--C10 of
Definition~\ref{def:catalogue-entries} at a parameter tuple in its declared
domain. Then:
\begin{enumerate}[label=(\roman*),leftmargin=2.2em]
\item $\Gmat_\Sigma^\dagger\Gmat_\Sigma\preceq\id_\Sigma$, so $\Gmat_\Sigma$
  is a \tref{local element} in the sense of Definition~\ref{def:embed};
\item $\Gmat_\Sigma$ is lossless,
  $\Gmat_\Sigma^\dagger\Gmat_\Sigma=\id_\Sigma$, exactly on the subfamily
  named in Table~\ref{tab:catalogue};
\item for C1--C9, a matrix whose parameter tuple passes the domain comparison
  of (E1) (Definition~\ref{def:graph}) is contractive by (i); C10 is the only
  entry whose declared domain is the contractivity condition itself, so that
  its check of Theorem~\ref{thm:compilation-well-defined}\,(iv)(E1) is a
  spectral test on the supplied matrix.
\end{enumerate}
\end{theorem}

\begin{proof}
\emph{Idea.} Each of the entries C1--C9 is exhibited in the form $D_2VD_1$ of
Lemma~\ref{lem:attenuated}, after which only the unitarity of the core has to
be checked, and the lemma then gives (i) and the lossless case (ii) together.
Entry C10 is decided by its hypothesis, checked on the supplied matrix.

\emph{C1.} $\Bsp(x,\varphi)$ is unitary: its first column has squared norm
$x+(1-x)=1$, its second $(1-x)+x=1$, and their inner product is
$\sqrt{x}\sqrt{1-x}e^{i\varphi}+\overline{\sqrt{1-x}e^{-i\varphi}}(-\sqrt{x})
=\sqrt{x(1-x)}e^{i\varphi}-\sqrt{x(1-x)}e^{i\varphi}=0$. Tensoring with
$\id_2$ preserves unitarity, and
$\operatorname{diag}(\sqrt{1-y_1},\sqrt{1-y_2})\otimes\id_2$ is a positive
semidefinite contraction because $1-y_i\in[0,1]$. Lemma~\ref{lem:attenuated}
with $D_1=\id$ gives admissibility and losslessness exactly at $y_1=y_2=0$.
(At $\varphi=0$ the core is real with $\det\Bsp=-1$; that sign is an artefact
of the convention, see Remark~\ref{rem:bsconv}.)

\emph{C2.} $\Bsp(t_H,\varphi_H)\oplus\Bsp(t_V,\varphi_V)$ is a direct sum of
unitaries, hence unitary; conjugating by the permutation $\Xi$ preserves
unitarity; the prefactor is again a positive semidefinite diagonal contraction.
The polarization-dependent beamsplitter of entry C1 is this matrix written in
the reordered basis $(p_1H,p_2H,p_1V,p_2V)$, so the same argument covers it.
At $t_H=1$, $t_V=0$, $\varphi_H=\varphi_V=0$, $y_1=y_2=0$ the element is the
signed permutation
$p_1H\mapsto p_1H$, $p_1V\mapsto p_2V$, $p_2H\mapsto -p_2H$,
$p_2V\mapsto p_1V$, which is the intended routing up to per-mode phases.

\emph{C3.} $\Rot(\theta)$ is real orthogonal ($\Rot(\theta)^{\transpose}
\Rot(\theta)=\id$ by $\cos^2+\sin^2=1$) and
$\operatorname{diag}(1,e^{i\delta})$ is unitary, so
$\Wplate(\delta,\theta)$ is a product of unitaries times the unimodular scalar
$e^{-i\delta/2}$, hence unitary; the expanded form in \eqref{eq:retarder} is
the product carried out, and \eqref{eq:hwpqwp} is $\delta=\pi$, using
$e^{i\pi}=-1$, $\cos^2\theta-\sin^2\theta=\cos2\theta$ and
$2\sin\theta\cos\theta=\sin2\theta$, with the scalar $e^{-i\pi/2}=-i$. The
prefactor $\sqrt{1-y}$ is a scalar contraction.

\emph{C4.} $\Rot(\rho)$ is real orthogonal, hence unitary, with no attenuation
factor.

\emph{C5.} $\operatorname{diag}(e^{i\varphi_H},e^{i\varphi_V})F$ is a product
of unitaries ($F$ is diagonal with entries of modulus one in both cases), and
$\operatorname{diag}(\sqrt{R_H},\sqrt{R_V})$ is a positive semidefinite
contraction, equal to $\id$ exactly when $R_H=R_V=1$.

\emph{C6, C9.} Diagonal with unimodular entries, respectively a permutation
matrix: both unitary, both with $D_1=D_2=\id$.

\emph{C7.} $V=\id$, $D_2=\operatorname{diag}(\sqrt{\eta_m})$, a positive
semidefinite contraction equal to $\id$ exactly when every $\eta_m=1$.

\emph{C8.} The two displayed forms agree because
$\Rot(\theta)\ket{H}=\ket{e_\theta}$ and $\Rot(\theta)\ket{V}=
\ket{e_\theta^{\perp}}$, so conjugating $\operatorname{diag}(\sqrt t,
\sqrt{t_\perp})$ by $\Rot(\theta)$ is the spectral decomposition written in
the rotated basis. The result is Hermitian positive semidefinite with
eigenvalues $\sqrt t,\sqrt{t_\perp}\in[0,1]$, so it is itself a positive
semidefinite contraction: Lemma~\ref{lem:attenuated} applies with
$V=\id$, $D_1=\id$, and equality holds exactly at $t=t_\perp=1$.

\emph{C10.} The condition is the hypothesis, checked on the supplied matrix;
losslessness is the same check with equality.

\emph{Item (iii).} For C1--C9 it is (i), which holds on the whole declared
domain; for C10 the declared domain is, by
Definition~\ref{def:catalogue-entries}, the set of blocks $\Gmat_\Sigma$ with
$\Gmat_\Sigma^\dagger\Gmat_\Sigma\preceq\id_\Sigma$.
\end{proof}
